\section{RELATED WORK AND MOTIVATION}
\label{sec:background}

%This section reviews the privacy risks of federated learning, the limitations of existing defences, and the technical constraints that motivate our design.


%% ─────────────────────────────────────────────────────
\subsection{Privacy Risks of Multi-Round Federated Learning}
\label{sec:bg_multiround}

Federated learning (FL)~\cite{mcmahan2017communication} retains raw data locally but repeatedly transmits model updates with each round, presenting an attack surface. Shared gradients can be inverted to reconstruct training images~\cite{zhu2019deep}, with recent attacks achieving exact reconstruction for batches up to $b \lesssim 25$ by exploiting ReLU-induced gradient sparsity~\cite{dimitrov2024spear}. Active and passive membership inference attacks exploit intermediate parameter updates~\cite{nasr2019comprehensive}, and GAN-based attacks reconstruct training data by exploiting evolving model dynamics across rounds~\cite{hitaj2017deep}. Even with secure aggregation, information accumulates: standard random user selection can leak individual user models within $O(N)$ rounds~\cite{so2023securing} N being the number of clients, and client-specific properties can be recovered from aggregated updates~\cite{kerkouche2023property}. Consequently, repeated communication is inherently risky regardless of raw data sharing policies~\cite{carletti2025sok}. HE-IFD eliminates this attack surface by restricting communication to a single encrypted round.


%% ─────────────────────────────────────────────────────
\subsection{Privacy Risks of Federated Distillation}
\label{sec:bg_server_inference}

Federated distillation is a class of FL methods in which clients train local teacher models and transfer knowledge to a server-side student by sharing model outputs such as
logits, soft labels, or intermediate features, rather than gradients. Several protocols utilize this paradigm: FedMD~\cite{li2019fedmd} and DS-FL~\cite{itahara2021dsfl} have clients share softmax predictions on a shared public dataset, FedDF~\cite{lin2020feddf} distils an ensemble of heterogeneous client models into a single server model, and Cronus~\cite{chang2019cronus} adds robust aggregation via spectral filtering of client predictions. However, plaintext logits and features leak private information even without gradient access, posing a central risk to unencrypted federated distillation. Model confidence values inherently provide sufficient information for inversion attacks, enabling adversaries to reconstruct training data features directly~\cite{fredrikson2015model}.

Against FedMD-style distillation, the Paired-Logits Inversion attack~\cite{takahashi2023breaching} allows a malicious server to recover high-fidelity private images by exploiting the confidence gap between its predictions and those of the clients. Even without reconstructing individual samples, an honest-but-curious server can infer a client's label distribution solely from aggregated predictions~\cite{gu2023ldia, shi2025unveiling}.

% Against FedMD-style distillation, the Paired-Logits Inversion attack~\cite{takahashi2023breaching} allows a malicious server to recover high-fidelity private images by exploiting the confidence gap between its predictions and those of the clients. Shi et al.~\cite{shi2025unveiling} demonstrate that client logit vectors encode the memorization of private data, enabling highly accurate label distribution and membership inference attacks against FedMD~\cite{li2019fedmd}, DS-FL~\cite{itahara2021dsfl}, and Cronus~\cite{chang2019cronus}. Additionally, an honest-but-curious server can infer a client's label distribution solely from aggregated predictions without reconstructing individual samples~\cite{gu2023ldia}.

Membership inference attacks pose a distinct threat by revealing whether a specific sample was used for training. FD-Leaks~\cite{yang2023fdleaks} uses a malicious client's local model as a shadow model, while GradDiff and GradDrift~\cite{wang2024graddiff} enable gradient-based inference by both servers and clients. More recent likelihood-ratio methods achieve stronger results: Co-operative LiRA and distillation-based LiRA variants~\cite{shi2025unveiling} attain high true-positive rates at low false-positive rates against FedMD and DS-FL. Distillation itself amplifies these vulnerabilities. As an example, students can exhibit up to $+9.81\%$ higher MIA susceptibility than their teachers~\cite{cui2025mia}, and the GLiRA framework~\cite{galichin2025glira} shows that knowledge-distillation-based shadow models improve likelihood-ratio attacks over the original LiRA baseline~\cite{carlini2022membership}. Even transmitting hard labels instead of soft labels only mitigates rather than eliminates this risk, as discrete predictions still leak label distributions~\cite{shao2024selective}.

% Membership inference attacks (MIA) further exacerbate these risks. FD-Leaks~\cite{yang2023fdleaks} uses a malicious client's local model as a shadow model, while GradDiff and GradDrift~\cite{wang2024graddiff} enable gradient-based inference by servers and clients. Distillation amplifies these vulnerabilities: students can exhibit up to $+9.81\%$ higher MIA susceptibility than their teachers~\cite{cui2025mia}, and the GLiRA framework~\cite{galichin2025glira} demonstrates that knowledge-distillation-based shadow models improve likelihood-ratio attacks over the original LiRA baseline~\cite{carlini2022membership}. Co-operative LiRA and distillation-based LiRA variants~\cite{shi2025unveiling} achieve high true-positive rates at low false-positive rates against FedMD- and DS-FL-style protocols. Even transmitting hard labels instead of soft labels only mitigates, rather than eliminates, this risk; discrete predictions still leak label distributions~\cite{shao2024selective}.

% \textbf{Beyond privacy: poisoning and model stealing.} 
% The attack surface extends beyond privacy violations. Logit poisoning attacks~\cite{tang2024logits} allow adversaries to manipulate shared predictions and corrupt the distilled student model. A semi-honest server can reconstruct client models from observed logits~\cite{wan2024privacy}, enabling model stealing. Even parametric feature sharing is unsafe: feature inversion can recover private images with SSIM values approaching $0.997$~\cite{beitollahi2024parametric}.

This diverse attack landscape demonstrates that plaintext knowledge transfer is fundamentally unsafe regardless of format. Only cryptographic protection eliminates these attack vectors entirely without degrading the distillation signal. HE-IFD encrypts all transferred knowledge under CKKS before upload, ensuring that the server never observes any plaintext supervision signal.



%% ─────────────────────────────────────────────────────
\subsection{One-Shot Federated Distillation}
\label{sec:bg_oneshot}

One-shot federated learning~\cite{guha2019oneshot, wang2025towards} restricts client-server communication to a single round and is the family of FL designs that share our problem statement. Among the plaintext distillation systems in this family, FedMD~\cite{li2019fedmd} and DS-FL~\cite{itahara2021dsfl} have clients transfer softmax predictions on a shared public dataset; FedDF~\cite{lin2020feddf} distils an ensemble of heterogeneous client models into a server model; Cronus~\cite{chang2019cronus} adds robust aggregation by spectrally filtering the client predictions; DENSE~\cite{zhang2022dense} uploads full plaintext models and trains a generator on the server; Co-Boosting~\cite{dai2024coboosting} iteratively co-trains data and ensemble; and FuseFL~\cite{tang2024fusefl} is architecturally closest to HE-IFD in that it also decomposes the student into blocks and trains them sequentially. None of these methods provide cryptographic protection: every client contribution reaches the server in plaintext and is therefore subject to the inversion, label-distribution, and membership-inference attacks reviewed in Section~\ref{sec:bg_server_inference}. The distinguishing axis of HE-IFD relative to this family is privacy: the supervision signal is a ciphertext under collectively held keys, and the server's view is computationally indistinguishable from random regardless of the architectural similarity to FuseFL.

The smaller subset of one-shot FL systems that do offer formal privacy guarantees rely uniformly on differential privacy. FedKT~\cite{li2021fedkt} achieves $(\varepsilon, 0)$-DP through noisy voting on a public dataset; FedMD-NFDP~\cite{sun2021fedmdnfdp} uses a noise-free sampling mechanism for $(\varepsilon, \delta)$-DP; FedDiff~\cite{feddiff2024} applies DP to diffusion-model training; and FedDM~\cite{xiong2023feddm} supports optional DP on synthetic data. The structural cost of this approach is that one-shot settings consume the entire privacy budget in a single round, with no across-round amplification, so accuracy collapses toward random guessing at meaningful privacy levels ($\varepsilon \leq 1$)~\cite{sun2021fedmdnfdp, gad2024communication}. The distinguishing axis of HE-IFD relative to this DP family is the type of guarantee: cryptographic and noise-free rather than statistical and noise-coupled to utility.


%% ─────────────────────────────────────────────────────
\subsection{HE-Compatible Neural Network Design}
\label{sec:bg_he_networks}

Running neural networks on encrypted data requires that all operations reduce to addition and multiplication over ciphertexts. Standard operations such as ReLU, batch normalisation, and softmax are non-polynomial and cannot be evaluated under encryption. Prior work on \emph{encrypted inference}, where a pre-trained model with known plaintext weights evaluates on encrypted inputs, replaces ReLU with learnable polynomial activations~\cite{baruch2022methodology, garimella2021sisyphus} and absorbs batch normalisation parameters into preceding layers~\cite{ibarrondo2021fhebn}. Recent centralised work has addressed the stability of training such polynomial networks: tighter ReLU approximations~\cite{agamennone2025polynomial}, boundary-loss and gradient-clipping schedules for deeper polynomial models~\cite{alhossain2025training}, and reboot-style block-independent training for polynomial MLPs~\cite{pirillo2025reboot}. \emph{Encrypted training}~\cite{sav2021poseidon} is substantially harder than inference: both weights and data are ciphertexts, every learnable operation is ciphertext-by-ciphertext (consuming multiplicative levels), and the backward pass roughly doubles the depth of the forward pass. All of the works above operate on a single centralised model; HE-IFD is the first to address stable polynomial-network training under fully encrypted, heterogeneous federated supervision, where teachers trained on different distributions produce target features with incompatible magnitude ranges.


%% ─────────────────────────────────────────────────────
\subsection{Encrypted Federated Learning Systems}
\label{sec:bg_he_fl}

The works that HE-IFD is most directly compared with in Section~\ref{sec:experiments} are encrypted federated learning systems. They share the cryptographic privacy goal but differ from HE-IFD on three axes: the unit that is encrypted (gradient updates vs.\ teacher features), the number of communication rounds, and where the multiplicative-depth budget is spent.

\par\noindent\textbf{POSEIDON~\cite{sav2021poseidon}.} POSEIDON is the canonical multiparty-CKKS system for federated training; it supports the same threshold trust model as HE-IFD and shares its underlying multiparty-CKKS construction~\cite{mouchet2021multiparty}. The unit that is encrypted is the per-round gradient: every client encrypts its gradient updates under the collective key, the server homomorphically aggregates them, and the result is broadcast back. The protocol is therefore inherently multi-round, and each round consumes the full forward+backward depth of the model under encryption, requiring bootstrapping to be invoked across rounds. HE-IFD differs on every axis: a single round, encryption of intermediate teacher features rather than gradients, and a depth budget that is confined to a single block of the student because each block is trained on fresh ciphertexts.

\par\noindent\textbf{BatchCrypt~\cite{zhang2020batchcrypt}.} BatchCrypt instantiates the same multi-round encrypted-gradient pattern with batched Paillier encryption rather than CKKS. The scheme is additively homomorphic, which is sufficient for FedAvg-style gradient aggregation but rules out the polynomial-network training that HE-IFD performs server-side. The relevant axis of comparison in our experiments is communication cost: BatchCrypt's per-round cost is dominated by encrypted gradients, scales linearly with the number of rounds, and gives the curve a slope similar to encrypted POSEIDON despite the different scheme.

\par\noindent\textbf{FedSHE~\cite{wei2025fedshe}.} A more recent CKKS-based federated learning system that introduces adaptive packing strategies to compress the per-round encrypted-gradient upload, reducing the constant in front of the multi-round communication curve. FedSHE remains structurally multi-round and encrypts gradient updates rather than supervision signals; HE-IFD is orthogonal to this line of work in that it eliminates the multi-round pattern entirely rather than optimising its constants. The two designs are complementary: an adaptive-packing scheme analogous to FedSHE's would also reduce the size of HE-IFD's single ciphertext upload.

In the same broader category, Vanilla FedAvg~\cite{mcmahan2017communication} is the unencrypted multi-round baseline used in the communication-cost comparison, and FedML-HE~\cite{jin2023fedml} and CURE~\cite{kanpak2024cure} are further multi-round encrypted-FL systems we cite for context but do not include as primary curves.



%% ─────────────────────────────────────────────────────
% Motivation paragraphs moved to introduction.tex §I-B Motivation (label sec:intro_motivation) per R3-3.

\section{Preliminaries: CKKS Homomorphic Encryption}
\label{sec:ckks_prelim}

The Cheon--Kim--Kim--Song (CKKS) scheme~\cite{cheon2017ckks} is a homomorphic encryption scheme that supports approximate arithmetic over vectors of floating point numbers. Unlike exact HE schemes, CKKS is designed for machine learning workloads where small numerical errors are acceptable.

\textbf{Slots and batching.} A CKKS ciphertext encrypts a vector of up to $\mathcal{N}/2$ real numbers simultaneously, where $\mathcal{N}$ is the polynomial modulus degree (a power of two). Each position in this vector is called a \textit{slot}. For example, with $\mathcal{N} = 8{,}192$, a single ciphertext holds 4{,}096 slots. This SIMD (single instruction, multiple data) packing means that one encrypted operation applies simultaneously to all slots, enabling efficient batched computation over large feature maps.

\textbf{Multiplicative depth.} Each ciphertext carries a finite \textit{multiplicative depth budget} which is the maximum number of ciphertext-by-ciphertext multiplications that can be performed before the noise floor exceeds the signal. Each such multiplication consumes one multiplicative \textit{level}. Addition and ciphertext-by-plaintext multiplication are budget-free. Once the entire budget is exhausted, i.e., multiplicative levels are consumed, a \textit{bootstrapping} operation can refresh the budget, but at substantial computational cost (${\sim}$20 seconds per ciphertext). HE-IFD minimises bootstrapping by training each block independently on fresh ciphertexts (Section~\ref{sec:phase2}), making the per-block depth to ${\sim}$12 levels. Bootstrapping is used only during the cascade refinement step that composes blocks into the full model.


\textbf{Security.} CKKS security relies on the Ring Learning with Errors (RLWE) hardness assumption~\cite{lyubashevsky2010ideal}. Under this assumption, a ciphertext is computationally indistinguishable from a random ring element (IND-CPA security): an adversary who observes only ciphertexts learns nothing about the encrypted values. All intermediate results of homomorphic computation, i.e., activations, gradients, loss values, remain ciphertexts and inherit this guarantee.

% \section{RELATED WORK AND MOTIVATION}
% \label{sec:background}

% %This section reviews the privacy risks of federated learning, the limitations of existing defences, and the technical constraints that motivate our design.


% %% ─────────────────────────────────────────────────────
% \subsection{Privacy Risks of Multi-Round Federated Learning}
% \label{sec:bg_multiround}

% Federated learning (FL)~\cite{mcmahan2017communication} retains raw data locally but repeatedly transmits model updates with each round, presenting an attack surface. Shared gradients can be inverted to reconstruct training images~\cite{zhu2019deep}, with recent attacks achieving exact reconstruction for batches up to $b \lesssim 25$ by exploiting ReLU-induced gradient sparsity~\cite{dimitrov2024spear}. Active and passive membership inference attacks exploit intermediate parameter updates~\cite{nasr2019comprehensive}, and GAN-based attacks reconstruct training data by exploiting evolving model dynamics across rounds~\cite{hitaj2017deep}. Even with secure aggregation, information accumulates: standard random user selection can leak individual user models within $O(N)$ rounds~\cite{so2023securing} N being the number of clients, and client-specific properties can be recovered from aggregated updates~\cite{kerkouche2023property}. Consequently, repeated communication is inherently risky regardless of raw data sharing policies~\cite{carletti2025sok}.


% %% ─────────────────────────────────────────────────────
% \subsection{Privacy Risks of Federated Distillation}
% \label{sec:bg_server_inference}

% Federated distillation is a class of \textbf{one-shot} federated learning methods in which clients train local teacher models and transfer knowledge to a server-side student by sharing model outputs---such as
% logits, soft labels, or intermediate features---rather than gradients. Several protocols utilize this paradigm: FedMD~\cite{li2019fedmd} and DS-FL~\cite{itahara2021dsfl} have clients share softmax predictions on a shared public dataset, FedDF~\cite{lin2020feddf} distils an ensemble of heterogeneous client models into a single server model, and Cronus~\cite{chang2019cronus} adds robust aggregation via spectral filtering of client predictions. Plaintext logits and features leak private information even without gradient access, posing a central risk to unencrypted federated distillation. Model confidence values inherently provide sufficient information for inversion attacks, enabling adversaries to reconstruct training data features directly~\cite{fredrikson2015model}.

% Against FedMD-style distillation, the Paired-Logits Inversion attack~\cite{takahashi2023breaching} allows a malicious server to recover high-fidelity private images by exploiting the confidence gap between its predictions and those of the clients. Even without reconstructing individual samples, an honest-but-curious server can infer a client's label distribution solely from aggregated predictions~\cite{gu2023ldia, shi2025unveiling}.

% % Against FedMD-style distillation, the Paired-Logits Inversion attack~\cite{takahashi2023breaching} allows a malicious server to recover high-fidelity private images by exploiting the confidence gap between its predictions and those of the clients. Shi et al.~\cite{shi2025unveiling} demonstrate that client logit vectors encode the memorization of private data, enabling highly accurate label distribution and membership inference attacks against FedMD~\cite{li2019fedmd}, DS-FL~\cite{itahara2021dsfl}, and Cronus~\cite{chang2019cronus}. Additionally, an honest-but-curious server can infer a client's label distribution solely from aggregated predictions without reconstructing individual samples~\cite{gu2023ldia}.

% Membership inference attacks pose a distinct threat by revealing whether a specific sample was used for training.FD-Leaks~\cite{yang2023fdleaks} uses a malicious client's local model as a shadow model, while GradDiff and GradDrift~\cite{wang2024graddiff} enable gradient-based inference by both servers and clients. More recent likelihood-ratio methods achieve  stronger results: Co-operative LiRA and distillation-based LiRA variants~\cite{shi2025unveiling} attain high true-positive rates at low false-positive rates against FedMD and DS-FL. Distillation itself amplifies these vulnerabilities---students can exhibit up to $+9.81\%$ higher MIA susceptibility than their teachers~\cite{cui2025mia}, and the GLiRA framework~\cite{galichin2025glira} shows that knowledge-distillation-based shadow models improve likelihood-ratio attacks over the original LiRA baseline~\cite{carlini2022membership}. Even transmitting hard labels instead of soft labels only mitigates rather than eliminates this risk, as discrete predictions still leak label distributions~\cite{shao2024selective}.

% % Membership inference attacks (MIA) further exacerbate these risks. FD-Leaks~\cite{yang2023fdleaks} uses a malicious client's local model as a shadow model, while GradDiff and GradDrift~\cite{wang2024graddiff} enable gradient-based inference by servers and clients. Distillation amplifies these vulnerabilities: students can exhibit up to $+9.81\%$ higher MIA susceptibility than their teachers~\cite{cui2025mia}, and the GLiRA framework~\cite{galichin2025glira} demonstrates that knowledge-distillation-based shadow models improve likelihood-ratio attacks over the original LiRA baseline~\cite{carlini2022membership}. Co-operative LiRA and distillation-based LiRA variants~\cite{shi2025unveiling} achieve high true-positive rates at low false-positive rates against FedMD- and DS-FL-style protocols. Even transmitting hard labels instead of soft labels only mitigates, rather than eliminates, this risk; discrete predictions still leak label distributions~\cite{shao2024selective}.

% % \textbf{Beyond privacy: poisoning and model stealing.} 
% % The attack surface extends beyond privacy violations. Logit poisoning attacks~\cite{tang2024logits} allow adversaries to manipulate shared predictions and corrupt the distilled student model. A semi-honest server can reconstruct client models from observed logits~\cite{wan2024privacy}, enabling model stealing. Even parametric feature sharing is unsafe: feature inversion can recover private images with SSIM values approaching $0.997$~\cite{beitollahi2024parametric}.

% This diverse attack landscape demonstrates that plaintext knowledge transfer is fundamentally unsafe regardless of format. Only cryptographic protection eliminates these attack vectors entirely without degrading the distillation signal.


% %% ─────────────────────────────────────────────────────
% \subsection{HE-Compatible Neural Network Design}
% \label{sec:bg_he_networks}

% Encrypted distillation refers to the process of training a student model entirely on ciphertext inputs and targets, such that the server performing the distillation never observes any plaintext data or model outputs. This requires a student architecture in which all operations reduce to addition and multiplication over ciphertexts. Standard operations such as ReLU, batch normalisation, and softmax are non-polynomial and cannot be evaluated under encryption. Prior work addresses this for centralised inference by replacing ReLU with learnable polynomial activations~\cite{baruch2022methodology} and absorbing batch normalisation parameters into preceding layers at inference time~\cite{ibarrondo2021fhebn}. But these solutions focus on inference and do not address the federated distillation setting where the supervision signal itself is an encrypted ciphertext. A further challenge in multi-client settings is that heterogeneous client distributions produce features with incompatible magnitude ranges, making centralised normalisation inapplicable. Our work is the first to address stable training of such architectures under fully encrypted federated supervision.



% %% ─────────────────────────────────────────────────────
% \par\noindent\textbf{Motivation for One-Shot Communication and Homomorphic Encryption.}
% % \label{sec:bg_motivation}
% One-shot federated learning~\cite{guha2019oneshot, wang2025towards} limits client--server communication to a single round, eliminating the iterative exposure that facilitates multi-round attacks. Knowledge distillation suits this setting naturally: clients train local teachers and transfer knowledge to a server-side student in a single step. However, data heterogeneity causes inconsistent teacher logits, and as shown above, even this single round of plaintext transfer remains exploitable by the server. Encrypting the transferred knowledge is therefore necessary, not optional.

% Homomorphic encryption (HE)~\cite{cheon2017ckks} provides end-to-end cryptographic protection without injecting noise, avoiding the utility--privacy tradeoff of DP. However, applying HE to standard iterative training is prohibitively expensive: systems like BatchCrypt~\cite{zhang2020batchcrypt} remain significantly slower than plaintext methods, and training deep networks over many rounds rapidly exhausts the finite multiplicative depth budget in the Cheon--Kim--Kim--Song (CKKS) scheme~\cite{cheon2017ckks} where we furhther discuss in \ref{sec:ckks_prelim}. Bootstrapping~\cite{cheon2018bootstrapping} refreshes this budget but introduces substantial latency, rendering layer-wise application impractical. One-shot settings resolve this tension: ciphertexts are used once, and block-independent training with fresh ciphertexts confines the depth budget to a single block rather than the full network. This principle was recently validated for encrypted MLP training~\cite{pirillo2025reboot}. Our framework extends it to federated one-shot distillation over heterogeneous clients.

% DP partially addresses plaintext leakage but degrades the distillation signal. Applying local DP causes severe accuracy drops (up to $70\%$ on CIFAR-100 at $\epsilon = 0.1$~\cite{gad2024communication}), and one-shot settings consume the entire privacy budget immediately, precluding amplification across rounds. At meaningful privacy levels ($\epsilon \leq 1$), DP reduces FedMD to near random-guessing accuracy~\cite{sun2021fedmdnfdp}. HE avoids this tradeoff entirely.

% \subsection{Preliminaries: CKKS Homomorphic Encryption}
% \label{sec:ckks_prelim}

% The Cheon--Kim--Kim--Song (CKKS) scheme~\cite{cheon2017ckks} is a homomorphic encryption scheme that supports approximate arithmetic over vectors of floating point numbers. Unlike exact HE schemes, CKKS is designed for machine learning workloads where small numerical errors are acceptable.

% \textbf{Slots and batching.} A CKKS ciphertext encrypts a vector of up to $\mathcal{N}/2$ real numbers simultaneously, where $\mathcal{N}$ is the polynomial modulus degree (a power of two). Each position in this vector is called a \textit{slot}. For example, with $\mathcal{N} = 8{,}192$, a single ciphertext holds 4{,}096 slots. This SIMD (single instruction, multiple data) packing means that one encrypted operation applies simultaneously to all slots, enabling efficient batched computation over large feature maps.

% \textbf{Multiplicative depth.} Each ciphertext carries a finite \textit{multiplicative depth budget}: the maximum number of ciphertext-by-ciphertext multiplications that can be performed before the noise floor exceeds the signal. Each such multiplication consumes one multiplicative \textit{level}. Addition and ciphertext-by-plaintext multiplication are budget-free. Once all budget is exhausted ie. multiplicative levels are consumed, a \textit{bootstrapping} operation can refresh the budget, but at substantial computational cost. HE-IFD avoids bootstrapping entirely (Section~\ref{sec:phase2}) by supplying fresh ciphertexts at each block boundary.

% \textbf{Security.} CKKS security rests on the Ring Learning with Errors (RLWE) hardness assumption~\cite{lyubashevsky2010ideal}. Under this assumption, a ciphertext is computationally indistinguishable from a random ring element (IND-CPA security): an adversary who observes only ciphertexts learns nothing about the encrypted values. All intermediate results of homomorphic computation, i.e., activations, gradients, loss values, remain ciphertexts and inherit this guarantee.


% \section{RELATED WORK AND MOTIVATION}
% \label{sec:background}

% %This section reviews the privacy risks of federated learning, the limitations of existing defences, and the technical constraints that motivate our design.


% %% ─────────────────────────────────────────────────────
% \subsection{Privacy Risks of Multi-Round Federated Learning}
% \label{sec:bg_multiround}

% Federated learning (FL)~\cite{mcmahan2017communication} retains raw data locally but repeatedly transmits model updates with each round, presenting an attack surface. Shared gradients can be inverted to reconstruct training images~\cite{zhu2019deep}, with recent attacks achieving exact reconstruction for batches up to $b \lesssim 25$ by exploiting ReLU-induced gradient sparsity~\cite{dimitrov2024spear}. Active and passive membership inference attacks exploit intermediate parameter updates~\cite{nasr2019comprehensive}, and GAN-based attacks reconstruct training data by exploiting evolving model dynamics across rounds~\cite{hitaj2017deep}. Even with secure aggregation, information accumulates: standard random user selection can leak individual user models within $O(N)$ rounds~\cite{so2023securing} N being the number of clients, and client-specific properties can be recovered from aggregated updates~\cite{kerkouche2023property}. Consequently, repeated communication is inherently risky regardless of raw data sharing policies~\cite{carletti2025sok}.


% %% ─────────────────────────────────────────────────────
% \subsection{Privacy Risks of Federated Distillation}
% \label{sec:bg_server_inference}

% Federated distillation is a class of \textbf{one-shot} federated learning methods in which clients train local teacher models and transfer knowledge to a server-side student by sharing model outputs---such as
% logits, soft labels, or intermediate features---rather than gradients. Several protocols utilize this paradigm: FedMD~\cite{li2019fedmd} and DS-FL~\cite{itahara2021dsfl} have clients share softmax predictions on a shared public dataset, FedDF~\cite{lin2020feddf} distils an ensemble of heterogeneous client models into a single server model, and Cronus~\cite{chang2019cronus} adds robust aggregation via spectral filtering of client predictions. Plaintext logits and features leak private information even without gradient access, posing a central risk to unencrypted federated distillation. Model confidence values inherently provide sufficient information for inversion attacks, enabling adversaries to reconstruct training data features directly~\cite{fredrikson2015model}.

% Against FedMD-style distillation, the Paired-Logits Inversion attack~\cite{takahashi2023breaching} allows a malicious server to recover high-fidelity private images by exploiting the confidence gap between its predictions and those of the clients. Even without reconstructing individual samples, an honest-but-curious server can infer a client's label distribution solely from aggregated predictions~\cite{gu2023ldia, shi2025unveiling}.

% % Against FedMD-style distillation, the Paired-Logits Inversion attack~\cite{takahashi2023breaching} allows a malicious server to recover high-fidelity private images by exploiting the confidence gap between its predictions and those of the clients. Shi et al.~\cite{shi2025unveiling} demonstrate that client logit vectors encode the memorization of private data, enabling highly accurate label distribution and membership inference attacks against FedMD~\cite{li2019fedmd}, DS-FL~\cite{itahara2021dsfl}, and Cronus~\cite{chang2019cronus}. Additionally, an honest-but-curious server can infer a client's label distribution solely from aggregated predictions without reconstructing individual samples~\cite{gu2023ldia}.

% Membership inference attacks pose a distinct threat by revealing whether a specific sample was used for training.FD-Leaks~\cite{yang2023fdleaks} uses a malicious client's local model as a shadow model, while GradDiff and GradDrift~\cite{wang2024graddiff} enable gradient-based inference by both servers and clients. More recent likelihood-ratio methods achieve  stronger results: Co-operative LiRA and distillation-based LiRA variants~\cite{shi2025unveiling} attain high true-positive rates at low false-positive rates against FedMD and DS-FL. Distillation itself amplifies these vulnerabilities---students can exhibit up to $+9.81\%$ higher MIA susceptibility than their teachers~\cite{cui2025mia}, and the GLiRA framework~\cite{galichin2025glira} shows that knowledge-distillation-based shadow models improve likelihood-ratio attacks over the original LiRA baseline~\cite{carlini2022membership}. Even transmitting hard labels instead of soft labels only mitigates rather than eliminates this risk, as discrete predictions still leak label distributions~\cite{shao2024selective}.

% % Membership inference attacks (MIA) further exacerbate these risks. FD-Leaks~\cite{yang2023fdleaks} uses a malicious client's local model as a shadow model, while GradDiff and GradDrift~\cite{wang2024graddiff} enable gradient-based inference by servers and clients. Distillation amplifies these vulnerabilities: students can exhibit up to $+9.81\%$ higher MIA susceptibility than their teachers~\cite{cui2025mia}, and the GLiRA framework~\cite{galichin2025glira} demonstrates that knowledge-distillation-based shadow models improve likelihood-ratio attacks over the original LiRA baseline~\cite{carlini2022membership}. Co-operative LiRA and distillation-based LiRA variants~\cite{shi2025unveiling} achieve high true-positive rates at low false-positive rates against FedMD- and DS-FL-style protocols. Even transmitting hard labels instead of soft labels only mitigates, rather than eliminates, this risk; discrete predictions still leak label distributions~\cite{shao2024selective}.

% % \textbf{Beyond privacy: poisoning and model stealing.} 
% % The attack surface extends beyond privacy violations. Logit poisoning attacks~\cite{tang2024logits} allow adversaries to manipulate shared predictions and corrupt the distilled student model. A semi-honest server can reconstruct client models from observed logits~\cite{wan2024privacy}, enabling model stealing. Even parametric feature sharing is unsafe: feature inversion can recover private images with SSIM values approaching $0.997$~\cite{beitollahi2024parametric}.

% This diverse attack landscape demonstrates that plaintext knowledge transfer is fundamentally unsafe regardless of format. Only cryptographic protection eliminates these attack vectors entirely without degrading the distillation signal.


% %% ─────────────────────────────────────────────────────
% \subsection{HE-Compatible Neural Network Design}
% \label{sec:bg_he_networks}

% Encrypted distillation refers to the process of training a student model entirely on ciphertext inputs and targets, such that the server performing the distillation never observes any plaintext data or model outputs. This requires a student architecture in which all operations reduce to addition and multiplication over ciphertexts. Standard operations such as ReLU, batch normalisation, and softmax are non-polynomial and cannot be evaluated under encryption. Prior work addresses this for centralised inference by replacing ReLU with learnable polynomial activations~\cite{baruch2022methodology} and absorbing batch normalisation parameters into preceding layers at inference time~\cite{ibarrondo2021fhebn}. But these solutions focus on inference and do not address the federated distillation setting where the supervision signal itself is an encrypted ciphertext. A further challenge in multi-client settings is that heterogeneous client distributions produce features with incompatible magnitude ranges, making centralised normalisation inapplicable. Our work is the first to address stable training of such architectures under fully encrypted federated supervision.



% %% ─────────────────────────────────────────────────────
% \par\noindent\textbf{Motivation for One-Shot Communication and Homomorphic Encryption.}
% % \label{sec:bg_motivation}
% One-shot federated learning~\cite{guha2019oneshot, wang2025towards} limits client--server communication to a single round, eliminating the iterative exposure that facilitates multi-round attacks. Knowledge distillation suits this setting naturally: clients train local teachers and transfer knowledge to a server-side student in a single step. However, data heterogeneity causes inconsistent teacher logits, and as shown above, even this single round of plaintext transfer remains exploitable by the server. Encrypting the transferred knowledge is therefore necessary, not optional.

% Homomorphic encryption (HE)~\cite{cheon2017ckks} provides end-to-end cryptographic protection without injecting noise, avoiding the utility--privacy tradeoff of DP. However, applying HE to standard iterative training is prohibitively expensive: systems like BatchCrypt~\cite{zhang2020batchcrypt} remain significantly slower than plaintext methods, and training deep networks over many rounds rapidly exhausts the finite multiplicative depth budget in the Cheon--Kim--Kim--Song (CKKS) scheme~\cite{cheon2017ckks} where we furhther discuss in \ref{sec:ckks_prelim}. Bootstrapping~\cite{cheon2018bootstrapping} refreshes this budget but introduces substantial latency, rendering layer-wise application impractical. One-shot settings resolve this tension: ciphertexts are used once, and block-independent training with fresh ciphertexts confines the depth budget to a single block rather than the full network. This principle was recently validated for encrypted MLP training~\cite{pirillo2025reboot}. Our framework extends it to federated one-shot distillation over heterogeneous clients.

% DP partially addresses plaintext leakage but degrades the distillation signal. Applying local DP causes severe accuracy drops (up to $70\%$ on CIFAR-100 at $\epsilon = 0.1$~\cite{gad2024communication}), and one-shot settings consume the entire privacy budget immediately, precluding amplification across rounds. At meaningful privacy levels ($\epsilon \leq 1$), DP reduces FedMD to near random-guessing accuracy~\cite{sun2021fedmdnfdp}. HE avoids this tradeoff entirely.

% \subsection{Preliminaries: CKKS Homomorphic Encryption}
% \label{sec:ckks_prelim}

% The Cheon--Kim--Kim--Song (CKKS) scheme~\cite{cheon2017ckks} is a homomorphic encryption scheme that supports approximate arithmetic over vectors of floating point numbers. Unlike exact HE schemes, CKKS is designed for machine learning workloads where small numerical errors are acceptable.

% \textbf{Slots and batching.} A CKKS ciphertext encrypts a vector of up to $\mathcal{N}/2$ real numbers simultaneously, where $\mathcal{N}$ is the polynomial modulus degree (a power of two). Each position in this vector is called a \textit{slot}. For example, with $\mathcal{N} = 8{,}192$, a single ciphertext holds 4{,}096 slots. This SIMD (single instruction, multiple data) packing means that one encrypted operation applies simultaneously to all slots, enabling efficient batched computation over large feature maps.

% \textbf{Multiplicative depth.} Each ciphertext carries a finite \textit{multiplicative depth budget}: the maximum number of ciphertext-by-ciphertext multiplications that can be performed before the noise floor exceeds the signal. Each such multiplication consumes one multiplicative \textit{level}. Addition and ciphertext-by-plaintext multiplication are budget-free. Once all budget is exhausted ie. multiplicative levels are consumed, a \textit{bootstrapping} operation can refresh the budget, but at substantial computational cost. HE-IFD avoids bootstrapping entirely (Section~\ref{sec:phase2}) by supplying fresh ciphertexts at each block boundary.

% \textbf{Security.} CKKS security rests on the Ring Learning with Errors (RLWE) hardness assumption~\cite{lyubashevsky2010ideal}. Under this assumption, a ciphertext is computationally indistinguishable from a random ring element (IND-CPA security): an adversary who observes only ciphertexts learns nothing about the encrypted values. All intermediate results of homomorphic computation, i.e., activations, gradients, loss values, remain ciphertexts and inherit this guarantee.

% \section{RELATED WORK AND MOTIVATION}
% \label{sec:background}

% %This section reviews the privacy risks of federated learning, the limitations of existing defences, and the technical constraints that motivate our design.


% %% ─────────────────────────────────────────────────────
% \subsection{Privacy Risks of Multi-Round Federated Learning}
% \label{sec:bg_multiround}

% Federated learning (FL)~\cite{mcmahan2017communication} retains raw data locally but repeatedly transmits model updates with each round, presenting an attack surface. Shared gradients can be inverted to reconstruct training images~\cite{zhu2019deep}, with recent attacks achieving exact reconstruction for batches up to $b \lesssim 25$ by exploiting ReLU-induced gradient sparsity~\cite{dimitrov2024spear}. Active and passive membership inference attacks exploit intermediate parameter updates~\cite{nasr2019comprehensive}, and GAN-based attacks reconstruct training data by exploiting evolving model dynamics across rounds~\cite{hitaj2017deep}. Even with secure aggregation, information accumulates: standard random user selection can leak individual user models within $O(N)$ rounds~\cite{so2023securing} N being the number of clients, and client-specific properties can be recovered from aggregated updates~\cite{kerkouche2023property}. Consequently, repeated communication is inherently risky regardless of raw data sharing policies~\cite{carletti2025sok}.


% %% ─────────────────────────────────────────────────────
% \subsection{Privacy Risks of Federated Distillation}
% \label{sec:bg_server_inference}

% Federated distillation is a class of \textbf{one-shot} federated learning methods in which clients train local teacher models and transfer knowledge to a server-side student by sharing model outputs---such as
% logits, soft labels, or intermediate features---rather than gradients. Several protocols utilize this paradigm: FedMD~\cite{li2019fedmd} and DS-FL~\cite{itahara2021dsfl} have clients share softmax predictions on a shared public dataset, FedDF~\cite{lin2020feddf} distils an ensemble of heterogeneous client models into a single server model, and Cronus~\cite{chang2019cronus} adds robust aggregation via spectral filtering of client predictions. Plaintext logits and features leak private information even without gradient access, posing a central risk to unencrypted federated distillation. Model confidence values inherently provide sufficient information for inversion attacks, enabling adversaries to reconstruct training data features directly~\cite{fredrikson2015model}.

% Against FedMD-style distillation, the Paired-Logits Inversion attack~\cite{takahashi2023breaching} allows a malicious server to recover high-fidelity private images by exploiting the confidence gap between its predictions and those of the clients. Even without reconstructing individual samples, an honest-but-curious server can infer a client's label distribution solely from aggregated predictions~\cite{gu2023ldia, shi2025unveiling}.

% % Against FedMD-style distillation, the Paired-Logits Inversion attack~\cite{takahashi2023breaching} allows a malicious server to recover high-fidelity private images by exploiting the confidence gap between its predictions and those of the clients. Shi et al.~\cite{shi2025unveiling} demonstrate that client logit vectors encode the memorization of private data, enabling highly accurate label distribution and membership inference attacks against FedMD~\cite{li2019fedmd}, DS-FL~\cite{itahara2021dsfl}, and Cronus~\cite{chang2019cronus}. Additionally, an honest-but-curious server can infer a client's label distribution solely from aggregated predictions without reconstructing individual samples~\cite{gu2023ldia}.

% Membership inference attacks pose a distinct threat by revealing whether a specific sample was used for training.FD-Leaks~\cite{yang2023fdleaks} uses a malicious client's local model as a shadow model, while GradDiff and GradDrift~\cite{wang2024graddiff} enable gradient-based inference by both servers and clients. More recent likelihood-ratio methods achieve  stronger results: Co-operative LiRA and distillation-based LiRA variants~\cite{shi2025unveiling} attain high true-positive rates at low false-positive rates against FedMD and DS-FL. Distillation itself amplifies these vulnerabilities---students can exhibit up to $+9.81\%$ higher MIA susceptibility than their teachers~\cite{cui2025mia}, and the GLiRA framework~\cite{galichin2025glira} shows that knowledge-distillation-based shadow models improve likelihood-ratio attacks over the original LiRA baseline~\cite{carlini2022membership}. Even transmitting hard labels instead of soft labels only mitigates rather than eliminates this risk, as discrete predictions still leak label distributions~\cite{shao2024selective}.

% % Membership inference attacks (MIA) further exacerbate these risks. FD-Leaks~\cite{yang2023fdleaks} uses a malicious client's local model as a shadow model, while GradDiff and GradDrift~\cite{wang2024graddiff} enable gradient-based inference by servers and clients. Distillation amplifies these vulnerabilities: students can exhibit up to $+9.81\%$ higher MIA susceptibility than their teachers~\cite{cui2025mia}, and the GLiRA framework~\cite{galichin2025glira} demonstrates that knowledge-distillation-based shadow models improve likelihood-ratio attacks over the original LiRA baseline~\cite{carlini2022membership}. Co-operative LiRA and distillation-based LiRA variants~\cite{shi2025unveiling} achieve high true-positive rates at low false-positive rates against FedMD- and DS-FL-style protocols. Even transmitting hard labels instead of soft labels only mitigates, rather than eliminates, this risk; discrete predictions still leak label distributions~\cite{shao2024selective}.

% % \textbf{Beyond privacy: poisoning and model stealing.} 
% % The attack surface extends beyond privacy violations. Logit poisoning attacks~\cite{tang2024logits} allow adversaries to manipulate shared predictions and corrupt the distilled student model. A semi-honest server can reconstruct client models from observed logits~\cite{wan2024privacy}, enabling model stealing. Even parametric feature sharing is unsafe: feature inversion can recover private images with SSIM values approaching $0.997$~\cite{beitollahi2024parametric}.

% This diverse attack landscape demonstrates that plaintext knowledge transfer is fundamentally unsafe regardless of format. Only cryptographic protection eliminates these attack vectors entirely without degrading the distillation signal.


% %% ─────────────────────────────────────────────────────
% \subsection{HE-Compatible Neural Network Design}
% \label{sec:bg_he_networks}

% Encrypted distillation refers to the process of training a student model entirely on ciphertext inputs and targets, such that the server performing the distillation never observes any plaintext data or model outputs. This requires a student architecture in which all operations reduce to addition and multiplication over ciphertexts. Standard operations such as ReLU, batch normalisation, and softmax are non-polynomial and cannot be evaluated under encryption. Prior work addresses this for centralised inference by replacing ReLU with learnable polynomial activations~\cite{baruch2022methodology} and absorbing batch normalisation parameters into preceding layers at inference time~\cite{ibarrondo2021fhebn}. But these solutions focus on inference and do not address the federated distillation setting where the supervision signal itself is an encrypted ciphertext. A further challenge in multi-client settings is that heterogeneous client distributions produce features with incompatible magnitude ranges, making centralised normalisation inapplicable. Our work is the first to address stable training of such architectures under fully encrypted federated supervision.



% %% ─────────────────────────────────────────────────────
% \par\noindent\textbf{Motivation for One-Shot Communication and Homomorphic Encryption.}
% % \label{sec:bg_motivation}
% One-shot federated learning~\cite{guha2019oneshot, wang2025towards} limits client--server communication to a single round, eliminating the iterative exposure that facilitates multi-round attacks. Knowledge distillation suits this setting naturally: clients train local teachers and transfer knowledge to a server-side student in a single step. However, data heterogeneity causes inconsistent teacher logits, and as shown above, even this single round of plaintext transfer remains exploitable by the server. Encrypting the transferred knowledge is therefore necessary, not optional.

% Homomorphic encryption (HE)~\cite{cheon2017ckks} provides end-to-end cryptographic protection without injecting noise, avoiding the utility--privacy tradeoff of DP. However, applying HE to standard iterative training is prohibitively expensive: systems like BatchCrypt~\cite{zhang2020batchcrypt} remain significantly slower than plaintext methods, and training deep networks over many rounds rapidly exhausts the finite multiplicative depth budget in the Cheon--Kim--Kim--Song (CKKS) scheme~\cite{cheon2017ckks} where we furhther discuss in \ref{sec:ckks_prelim}. Bootstrapping~\cite{cheon2018bootstrapping} refreshes this budget but introduces substantial latency, rendering layer-wise application impractical. One-shot settings resolve this tension: ciphertexts are used once, and block-independent training with fresh ciphertexts confines the depth budget to a single block rather than the full network. This principle was recently validated for encrypted MLP training~\cite{pirillo2025reboot}. Our framework extends it to federated one-shot distillation over heterogeneous clients.

% DP partially addresses plaintext leakage but degrades the distillation signal. Applying local DP causes severe accuracy drops (up to $70\%$ on CIFAR-100 at $\epsilon = 0.1$~\cite{gad2024communication}), and one-shot settings consume the entire privacy budget immediately, precluding amplification across rounds. At meaningful privacy levels ($\epsilon \leq 1$), DP reduces FedMD to near random-guessing accuracy~\cite{sun2021fedmdnfdp}. HE avoids this tradeoff entirely.

% \section{Preliminaries: CKKS Homomorphic Encryption}
% \label{sec:ckks_prelim}

% The Cheon--Kim--Kim--Song (CKKS) scheme~\cite{cheon2017ckks} is a homomorphic encryption scheme that supports approximate arithmetic over vectors of floating point numbers. Unlike exact HE schemes, CKKS is designed for machine learning workloads where small numerical errors are acceptable.

% \textbf{Slots and batching.} A CKKS ciphertext encrypts a vector of up to $\mathcal{N}/2$ real numbers simultaneously, where $\mathcal{N}$ is the polynomial modulus degree (a power of two). Each position in this vector is called a \textit{slot}. For example, with $\mathcal{N} = 8{,}192$, a single ciphertext holds 4{,}096 slots. This SIMD (single instruction, multiple data) packing means that one encrypted operation applies simultaneously to all slots, enabling efficient batched computation over large feature maps.

% \textbf{Multiplicative depth.} Each ciphertext carries a finite \textit{multiplicative depth budget}: the maximum number of ciphertext-by-ciphertext multiplications that can be performed before the noise floor exceeds the signal. Each such multiplication consumes one multiplicative \textit{level}. Addition and ciphertext-by-plaintext multiplication are budget-free. Once all budget is exhausted ie. multiplicative levels are consumed, a \textit{bootstrapping} operation can refresh the budget, but at substantial computational cost. HE-IFD avoids bootstrapping entirely as we discuss in \ref{sec:phase2_server} by supplying fresh ciphertexts at each block boundary.

% \par\noindent\textbf{Security.} CKKS security rests on the Ring Learning with Errors (RLWE) hardness assumption~\cite{lyubashevsky2010ideal}. Under this assumption, a ciphertext is computationally indistinguishable from a random ring element (IND-CPA security): an adversary who observes only ciphertexts learns nothing about the encrypted values. All intermediate results of homomorphic computation, i.e., activations, gradients, loss values, remain ciphertexts and inherit this guarantee.

